\documentclass[aspectratio=169]{beamer}
% Préambule Beamer commun aux huit séquences.
% Les decks restent volontairement compatibles avec Beamer standard.

\usepackage[T1]{fontenc}
\usepackage[utf8]{inputenc}
\usepackage[french]{babel}
\usepackage{graphicx}
\usepackage{booktabs}
\usepackage{hyperref}

\setbeamertemplate{navigation symbols}{}
\setbeamertemplate{footline}[frame number]

\newcommand{\decisionquestion}[1]{%
  \begin{block}{Question décisionnelle}
    #1
  \end{block}
}

\newcommand{\confidencelevel}[1]{\textbf{Niveau de confiance :} #1}

\graphicspath{{figures/}}
\title{Séquence 2 -- Comprendre la source}
\subtitle{Broker MQTT, topics, messages et complétude}
\author{École de l'air et de l'espace}
\date{}
\newcommand{\sourcefoot}[1]{\vfill\tiny\color{gray}Source : #1}
\begin{document}

\begin{frame}\titlepage\end{frame}

\begin{frame}{La connexion réussit ; la couverture reste à démontrer}
\large Un broker qui répond à une commande de connexion prouve seulement une disponibilité technique ponctuelle, à cet instant précis.\\[1em]
Il ne prouve pas que les cinq zones critiques de la base sont représentées dans ses messages, que ces messages sont récents, ni que les capteurs qui les ont produits fonctionnent encore.
\decisionquestion{L'état observé du broker suffit-il, à lui seul, pour soutenir une décision couvrant toute la base ?}
\end{frame}

\begin{frame}{À la fin, vous saurez borner une décision à la preuve disponible}
\begin{itemize}
\item distinguer précisément broker, topic, filtre, payload et retained, et savoir laquelle de ces notions prouve quoi ;
\item inventorier ce qui est effectivement observé, sans confondre observation et attendu ;
\item comparer cet inventaire à un référentiel attendu explicite pour mesurer une véritable complétude ;
\item qualifier une absence de message sans inventer sa cause probable ;
\item recommander une action, une confiance et une vérification proportionnées à ce périmètre observé.
\end{itemize}
\end{frame}

\begin{frame}{La séance suit quatre productions clairement identifiées}
\begin{enumerate}
\item \textbf{TP 1} -- explorer les topics disponibles sans surinterpréter ce qu'ils signifient ;
\item \textbf{TP 2} -- construire l'inventaire reproductible des messages observés ;
\item \textbf{TP 3} -- diagnostiquer la complétude en comparant l'attendu à l'observé ;
\item \textbf{TP 4} -- rédiger puis faire contester le brief décisionnel.
\end{enumerate}
Chaque TP se clôt de la même façon : une décision provisoire, une confiance, une preuve, une incertitude et une limite explicite.
\end{frame}

\section{Accroche}
\begin{frame}{Activité A -- Le broker répond : que décidez-vous ?}
Une relève doit confirmer la couverture thermique de cinq zones critiques avant 10 h. Un opérateur annonce seulement : « le broker répond et contient des messages retained ».
\begin{block}{Vote individuel}
A couverture suffisante \quad B inspection ciblée \quad C suspendre la conclusion \quad D impossible sans inventaire
\end{block}
Notez votre confiance de 0 à 100 \%, la preuve exacte dont vous disposez et l'information qui vous manque le plus pour trancher.
\end{frame}

\begin{frame}{Votre confiance doit mesurer la suffisance de la preuve}
\begin{columns}
\column{.48\textwidth}\textbf{Confiance utile}
\begin{itemize}
\item elle est explicitement reliée au périmètre observé, pas à l'ensemble de la base ;
\item elle s'appuie sur des traces vérifiables et citables ;
\item elle peut être révisée à la baisse sans que cela soit perçu comme un échec.
\end{itemize}
\column{.48\textwidth}\textbf{Fausse confiance}
\begin{itemize}
\item elle vient de l'aisance avec l'outil technique, pas de la preuve elle-même ;
\item elle s'exprime sans jamais préciser ce qui était attendu ;
\item elle confond le sentiment de certitude avec une certitude réellement démontrée.
\end{itemize}
\end{columns}
\end{frame}

\section{Modèle MQTT}
\begin{frame}{MQTT sépare publication, distribution et consommation}
\centering\Large Publisher $\rightarrow$ Broker $\rightarrow$ Subscriber\\[1em]
\normalsize Le publisher choisit librement un topic et construit un payload, sans savoir qui le lira ni combien d'abonnés existent. Le broker distribue ensuite ce message selon les abonnements actifs, sans en interpréter le sens. Le subscriber, enfin, est seul responsable d'interpréter ce contenu dans son propre contexte métier.
\sourcefoot{OASIS, MQTT Version 5.0, 2019.}
\end{frame}

\begin{frame}{Le broker sait router ; il ne connaît pas votre intention opérationnelle}
\begin{itemize}
\item Il associe les publications reçues aux abonnements correspondants, selon les règles du protocole.
\item Il peut conserver un dernier message retained par topic, pour le remettre à un futur abonné.
\item Il ne sait pas spontanément qu'une base aérienne doit compter cinq zones surveillées : cette connaissance métier n'existe nulle part dans le protocole.
\item Il ne certifie ni la calibration d'un capteur, ni l'exhaustivité métier d'un lot, ni la vérité de ce qui se passe réellement sur le terrain.
\end{itemize}
\end{frame}

\begin{frame}{Un topic est une adresse applicative, pas une preuve}
\small\texttt{airbase/batch002/battery-shelter-01/temperature}
\begin{itemize}
\item Chaque niveau de ce chemin est séparé par le caractère \texttt{/}, comme un chemin de fichier.
\item Cette hiérarchie est une convention choisie par les concepteurs du système, pas une règle imposée par MQTT lui-même.
\item Un nom cohérent, comme celui-ci, facilite grandement la construction d'un inventaire lisible.
\item Mais un nom plausible peut néanmoins être faux ou incomplet : rien n'empêche une configuration erronée de publier sous un topic qui semble correct sans l'être.
\end{itemize}
\end{frame}

\begin{frame}{Le filtre définit votre champ de vision}
\begin{tabular}{@{}p{.4\textwidth}p{.52\textwidth}@{}}
\toprule
Filtre & Ce qu'il peut couvrir \\
\midrule
\texttt{airbase/batch002/\#} & toute la branche batch002, quelle que soit la zone \\[2pt]
\texttt{airbase/+/+/temperature} & les mesures de température, trois niveaux après airbase \\[2pt]
\texttt{airbase/batch001/\#} & le mauvais lot -- rien de batch002 ne sera reçu \\
\bottomrule
\end{tabular}
\begin{alertblock}{Risque}
Un filtre trop étroit, ou pointant sur le mauvais lot, produit une absence apparente de messages sans qu'aucun capteur ne soit réellement en panne : le silence peut venir de notre abonnement, pas du terrain.
\end{alertblock}
\sourcefoot{OASIS MQTT 5.0, Topic Names and Topic Filters.}
\end{frame}

\begin{frame}{Prédiction -- que recevra cet abonnement ?}
\large\texttt{airbase/batch002/\#}
\begin{enumerate}
\item tous les messages jamais publiés sur ce lot, depuis le début ?
\item uniquement les publications futures correspondant à ce filtre ?
\item les messages retained déjà stockés au moment de l'abonnement, plus les publications futures ?
\item la liste des capteurs qui devraient exister sur la base ?
\end{enumerate}
Justifiez chaque réponse par le mécanisme du protocole, pas par intuition : nous vérifierons ensemble dans quelques minutes, à l'exécution.
\end{frame}

\begin{frame}{Le payload transporte la mesure et son contexte déclaré}
\small
\begin{tabular}{ll}
\toprule
Champ & Question utile \\
\midrule
\texttt{zone} & où la mesure prétend-elle s'appliquer ? \\
\texttt{sensor} & quelle grandeur physique est annoncée ? \\
\texttt{measured\_at} & quand le capteur déclare-t-il avoir mesuré ? \\
\texttt{message\_id} & peut-on retrouver cette publication précise ? \\
\texttt{value}, \texttt{unit} & quelle valeur est déclarée, dans quelle unité ? \\
\bottomrule
\end{tabular}
\begin{exampleblock}{Exemple réel du lot batch002}
Pour \texttt{battery-shelter-01} : \texttt{value=34.9}, \texttt{unit=celsius}, \texttt{measured\_at=2026-10-19T09:25:00Z}, \texttt{message\_id=battery-shelter-01-latest}. Chacun de ces cinq champs répond à l'une des questions ci-dessus ; aucun ne peut être omis sans perdre une partie de la portée de la mesure.
\end{exampleblock}
\end{frame}

\begin{frame}{Enveloppe et payload ne répondent pas aux mêmes questions}
\begin{columns}
\column{.48\textwidth}\textbf{Enveloppe d'extraction}
\begin{itemize}
\item le topic complet, qui décrit d'où vient le message ;
\item \texttt{received\_at}, l'instant où notre extraction l'a reçu ;
\item retained, qui indique s'il s'agissait d'un message déjà stocké.
\end{itemize}
\column{.48\textwidth}\textbf{Payload publié}
\begin{itemize}
\item la zone et le capteur déclarés par l'émetteur ;
\item \texttt{measured\_at}, l'instant déclaré de la mesure elle-même ;
\item la valeur mesurée et son unité.
\end{itemize}
\end{columns}
\begin{block}{Réflexe}
Nous conservons systématiquement les deux ensembles de champs : une donnée métier privée de son contexte de réception perd une partie essentielle de sa portée pour l'audit ultérieur.
\end{block}
\end{frame}

\begin{frame}{Retained signifie `` dernier état disponible '', pas `` état actuel garanti ''}
Lorsqu'un message est publié avec l'option retain, le broker en conserve une copie associée à ce topic. Tout nouvel abonné dont le filtre correspond à ce topic reçoit immédiatement cette copie, même s'il s'est abonné longtemps après la publication d'origine.
\begin{alertblock}{Trois non-conclusions}
Ce n'est pas un historique complet des mesures passées ; ce n'est pas une preuve que la mesure est fraîche ; ce n'est pas une preuve que le capteur qui l'a publiée fonctionne encore aujourd'hui.
\end{alertblock}
\sourcefoot{OASIS MQTT 5.0, Retain flag and retained messages ; Eclipse Mosquitto, mosquitto\_sub manual.}
\end{frame}

\begin{frame}{Vrai ou faux -- retained}
\begin{enumerate}
\item « Reçu maintenant » implique « mesuré maintenant ».
\item Un topic peut avoir au plus un message retained courant à un instant donné.
\item Un nouvel abonné peut recevoir un message retained immédiatement après son abonnement.
\item Le mécanisme retained fournit la liste des capteurs qui devraient exister.
\item Supprimer un message retained peut créer un silence observable sans qu'aucun capteur ne soit en panne.
\end{enumerate}
Pour chaque proposition, répondez vrai ou faux et donnez la conséquence décisionnelle si l'on se trompait.
\end{frame}

\begin{frame}{Auto-vérification avant de continuer}
\small
\begin{block}{Le broker connaît-il la liste des capteurs qui devraient exister sur la base ?}
Non : ni le protocole MQTT ni le broker ne possèdent cette connaissance métier. Elle vient toujours d'un référentiel externe que nous devrons fournir nous-mêmes.
\end{block}
\begin{block}{Que reçoit-on en s'abonnant à \texttt{airbase/batch002/\#} ?}
Les messages disponibles dont le topic commence par cette branche -- notamment les retained déjà stockés au moment de l'abonnement -- et non un historique complet de toutes les publications passées.
\end{block}
\begin{block}{Un message retained prouve-t-il que son capteur fonctionne encore ?}
Non : le broker peut continuer à servir le dernier état connu longtemps après l'arrêt réel du capteur qui l'a produit.
\end{block}
\end{frame}

\section{TP 1 -- Explorer}
\begin{frame}{TP 1 -- Votre objectif est d'observer, pas encore de conclure}
\begin{block}{Production attendue}
Une table à deux colonnes : « j'observe » / « je peux conclure », renseignée pour quatre messages réels.
\end{block}
Pour chaque message, relevez le topic complet, ses niveaux, la zone, le capteur, \texttt{measured\_at}, \texttt{received\_at} et retained. Consignez également le filtre exact utilisé pour l'extraction : sans lui, personne ne pourra vérifier la portée de votre observation.
\end{frame}

\begin{frame}[fragile]{TP 1 -- Extraire la branche batch002}
\small
\begin{verbatim}
export PYTHONPATH=src
python3 -m iot_decision.mqtt_tools extract \
  data/raw/batch002_observed.jsonl \
  --topic airbase/batch002/#
\end{verbatim}
Plan B si le broker n'est pas disponible : ouvrir directement \texttt{data/samples/batch002\_retained\_messages.jsonl}, qui contient exactement les mêmes messages.
\end{frame}

\begin{frame}[fragile]{TP 1 -- Vérifiez avant d'interpréter}
\small
\begin{verbatim}
test -f data/raw/batch002_observed.jsonl && echo "présent"
wc -l data/raw/batch002_observed.jsonl
head -n 1 data/raw/batch002_observed.jsonl
\end{verbatim}
\begin{block}{Sortie attendue}
\texttt{présent}, puis \texttt{4}, puis une enveloppe contenant topic, received\_at, retained et un payload complet.
\end{block}
\end{frame}

\begin{frame}{TP 1 -- Répondez avant de passer à l'inventaire}
\small
\begin{enumerate}
\item Quel filtre exact définit votre champ de vision pour cette extraction ?
\item Combien de topics avez-vous effectivement observés avec ce filtre ?
\item Quels champs viennent de l'enveloppe d'extraction ? Lesquels viennent du payload publié ?
\item Quelle horloge décrit l'instant de la mesure ? Laquelle décrit l'instant de la réception ?
\item Quelle phrase pouvez-vous défendre sans disposer d'un référentiel attendu ?
\end{enumerate}
\textbf{Trace :} une ligne « j'observe / je peux conclure » par message observé.
\end{frame}

\begin{frame}{TP 1 -- Contrôles avant interprétation}
\begin{enumerate}
\item Le fichier existe et n'est pas vide : sans ce contrôle, toute analyse porterait sur du vide.
\item Chaque ligne est un objet JSON valide contenant une enveloppe complète.
\item Le topic de chaque ligne commence bien par la branche demandée dans le filtre.
\item Retained et les deux horloges (\texttt{measured\_at}, \texttt{received\_at}) sont bien conservés dans chaque ligne.
\item L'effectif obtenu est une observation de cette extraction précise, jamais une preuve de complétude du lot entier.
\end{enumerate}
\end{frame}

\begin{frame}{Point d'étape -- quatre messages répondent à une question limitée}
Vous pouvez affirmer : « avec ce filtre et cette extraction, quatre topics ont été observés ».
\begin{alertblock}{Vous ne pouvez pas encore dire}
« Les quatre capteurs fonctionnent », « toutes les zones sont couvertes » ou « aucun message ne manque » : aucune de ces trois phrases n'est démontrée par le seul fait d'avoir reçu quatre messages.
\end{alertblock}
\end{frame}

\section{TP 2 -- Inventorier}
\begin{frame}{TP 2 -- Un inventaire conserve le lien entre message et contexte}
\begin{block}{Une ligne par topic observé}
topic, site\_id, zone, sensor, message\_id, measured\_at, received\_at, retained, payload\_fields -- neuf colonnes au total.
\end{block}
\begin{exampleblock}{Exemple réel}
La ligne pour \texttt{battery-shelter-01} porte le topic \texttt{airbase/batch002/battery-shelter-01/temperature}, \texttt{measured\_at=2026-10-19T09:25:00Z}, \texttt{received\_at=2026-10-19T09:30:00Z} et \texttt{retained=true}.
\end{exampleblock}
L'objectif de ce TP n'est pas de nettoyer les valeurs mesurées ; il est de décrire fidèlement le champ de preuve tel qu'il a été observé.
\end{frame}

\begin{frame}{Le topic complet empêche une agrégation prématurée}
Si vous ne conservez que le mot « température », vous perdez le lot, le site et la zone concernés. Si vous ne conservez que la zone, vous perdez la trace de la branche effectivement interrogée par votre filtre.
\decisionquestion{Quelle colonne de votre inventaire permet à un tiers de vérifier que votre travail porte bien sur le bon périmètre ?}
\end{frame}

\begin{frame}[fragile]{TP 2 -- Produire l'inventaire reproductible}
\scriptsize
\begin{verbatim}
python3 -m iot_decision.source_inventory_cli \
 data/samples/batch002_retained_messages.jsonl \
 data/samples/batch002_expected_sensors.csv \
 data/processed/batch002_inventory.csv \
 data/processed/batch002_completeness.json
\end{verbatim}
Ouvrez d'abord le fichier CSV pour contrôler l'inventaire brut. Laissez le diagnostic JSON fermé pour l'instant : nous y reviendrons après l'avoir construit manuellement.
\end{frame}

\begin{frame}[fragile]{TP 2 -- Contrôlez quatre lignes et quatre topics uniques}
\scriptsize
\begin{verbatim}
head -n 5 data/processed/batch002_inventory.csv
tail -n +2 data/processed/batch002_inventory.csv | wc -l
tail -n +2 data/processed/batch002_inventory.csv | cut -d, -f1 | sort -u | wc -l
\end{verbatim}
Les deux comptages doivent afficher \texttt{4} : quatre lignes, et quatre topics strictement uniques parmi ces lignes.
\end{frame}

\begin{frame}{TP 2 -- Les cinq réponses à remettre}
\small
\begin{enumerate}
\item Quelles quatre zones sont observées dans cet inventaire ?
\item Quelle colonne permet de retrouver précisément le périmètre interrogé par votre extraction ?
\item Pourquoi conserver \texttt{message\_id} et les deux horloges plutôt que la seule valeur mesurée ?
\item Que perdez-vous concrètement si vous regroupez immédiatement toutes les lignes par capteur ?
\item L'inventaire obtenu prouve-t-il le fonctionnement actuel des quatre capteurs correspondants ?
\end{enumerate}
\textbf{Fin du TP :} un CSV contrôlé, accompagné d'une mini-décision en six champs.
\end{frame}

\begin{frame}{TP 2 -- Revue croisée de l'inventaire}
Le binôme voisin vérifie méthodiquement :
\begin{itemize}
\item qu'il existe exactement une ligne par topic, et aucun doublon ;
\item que les quatre zones attendues à ce stade sont bien nommées explicitement ;
\item que les deux horloges et l'indicateur retained sont présents sur chaque ligne ;
\item que toutes les lignes appartiennent bien à la branche batch002 ;
\item qu'aucune affirmation du binôme ne dépasse strictement ce qui a été observé.
\end{itemize}
\end{frame}

\begin{frame}{Auto-vérification avant de continuer}
\small
\begin{block}{Pourquoi conserver le topic complet plutôt que la seule zone dans l'inventaire ?}
Parce que la zone seule ne permet plus de vérifier a posteriori quel filtre et quelle branche ont réellement été interrogés ; le topic complet, lui, le permet toujours.
\end{block}
\begin{block}{L'inventaire prouve-t-il que les capteurs observés fonctionnent réellement ?}
Non : il prouve seulement qu'un message correspondant à ce topic était disponible au moment de l'extraction, pas que le capteur continue de fonctionner aujourd'hui.
\end{block}
\begin{block}{Pourquoi parler de « topics observés » plutôt que de « capteurs fonctionnels » ?}
Parce que l'extraction constate la présence d'un message, sans tester ni l'alimentation, ni la calibration, ni le fonctionnement actuel du capteur qui l'a produit.
\end{block}
\end{frame}

\begin{frame}{Avant la pause -- formulez votre mini-décision}
\begin{block}{À remettre}
Une décision provisoire, une confiance, une preuve, une incertitude et une limite, en six champs courts.
\end{block}
Question à trancher : avez-vous, à ce stade, assez d'information pour employer l'expression « lot complet » ?
\end{frame}

\begin{frame}{Pause -- 10 minutes}
\centering\Large Au retour : complet par rapport à quoi, exactement ?
\end{frame}

\section{Complétude}
\begin{frame}{Un lot est un périmètre défini, pas un objet natif de MQTT}
Dans cette séance, le lot est borné explicitement par quatre éléments :
\begin{itemize}
\item le site concerné et la branche de topics \texttt{batch002} ;
\item l'instant et la durée de l'extraction réalisée ;
\item le filtre d'abonnement effectivement utilisé ;
\item la liste autorisée des topics attendus, fournie par un référentiel externe.
\end{itemize}
Changer un seul de ces quatre éléments change immédiatement la portée du diagnostic obtenu.
\end{frame}

\begin{frame}{La complétude exige un dénominateur explicite}
\centering\Large Couverture $= \frac{\text{topics attendus observés}}{\text{topics attendus}}$
\vspace{1em}
\normalsize Le numérateur vient directement de votre inventaire. Le dénominateur, lui, vient toujours d'un référentiel métier externe, qui doit avoir une autorité identifiée et une date de validité claire -- sans ce dénominateur, aucun pourcentage de couverture n'a de sens.
\end{frame}

\begin{frame}{Trois questions précèdent tout pourcentage}
\begin{enumerate}
\item Attendu par qui, précisément, et pour quelle mission ce référentiel a-t-il été établi ?
\item Ce référentiel reste-t-il valide à la date d'aujourd'hui, et pour ce site précis ?
\item La présence de chaque zone a-t-elle réellement la même importance pour la décision ?
\end{enumerate}
Un taux de couverture parfaitement exact peut malgré tout soutenir une mauvaise décision si son périmètre de départ est faux.
\end{frame}

\begin{frame}{Les métadonnées rendent le diagnostic contestable et vérifiable}
\begin{tabular}{ll}
\toprule
Métadonnée & Ce qu'elle permet de contester \\
\midrule
filtre utilisé & un champ de vision resté incomplet \\
instant de réception & une extraction lancée trop tôt ou trop courte \\
topic complet & une mauvaise branche ou une mauvaise zone \\
référentiel et sa version & un attendu devenu obsolète \\
retained & une confusion avec un flux de données vivant \\
\bottomrule
\end{tabular}
\vfill
\normalsize Sans ces métadonnées, un diagnostic de complétude ne pourrait ni être vérifié, ni être remis en cause par un tiers.
\end{frame}

\begin{frame}{Activité E -- Classez les affirmations}
Pour chacune, décidez si elle est directement observable, vérifiable seulement à l'aide d'un référentiel, ou non démontrable avec les moyens dont nous disposons ici :
\begin{itemize}
\item toutes les zones attendues sont présentes dans l'inventaire ;
\item aucun message n'a été perdu entre l'émission et l'extraction ;
\item chaque message retained observé est récent ;
\item les capteurs observés fonctionnent correctement en ce moment ;
\item le filtre utilisé a bien couvert toute la branche batch002.
\end{itemize}
\end{frame}

\begin{frame}{Auto-vérification avant de continuer}
\small
\begin{block}{Que faut-il toujours citer avant d'annoncer un taux de couverture ?}
Le référentiel exact utilisé comme dénominateur, son autorité, et sa date de validité pour ce site précis : un pourcentage sans ces trois éléments n'est pas vérifiable par un tiers.
\end{block}
\begin{block}{« Quatre topics uniques ont été inventoriés » est-elle une affirmation observable ?}
Oui, directement, par simple comptage dans le fichier d'inventaire produit.
\end{block}
\begin{block}{« Toutes les zones attendues sont présentes » est-elle, elle aussi, directement observable ?}
Non : elle n'est vérifiable qu'en la comparant à un référentiel attendu explicite, ce que nous allons faire maintenant dans le TP 3.
\end{block}
\end{frame}

\section{TP 3 -- Comparer}
\begin{frame}{TP 3 -- La matrice attendu/observé localise le manque}
\small
\begin{tabular}{llll}
\toprule
Zone attendue & Criticité & Observée ? & Preuve \\
\midrule
batteries & haute & ? & topic/ligne \\
transmissions & haute & ? & topic/ligne \\
informatique & haute & ? & topic/ligne \\
maintenance & moyenne & ? & topic/ligne \\
optronique & haute & ? & topic/ligne \\
\bottomrule
\end{tabular}
\vfill
\normalsize Complétez vous-même cette matrice, zone par zone, avant d'exécuter le diagnostic automatisé : c'est ce travail manuel qui vous permettra de contester le résultat produit par le script.
\end{frame}

\begin{frame}{TP 3 -- Travail avant l'automatisation}
\begin{enumerate}
\item ouvrez le référentiel attendu et lisez-en chaque ligne ;
\item cherchez, un par un, chaque topic attendu dans votre inventaire ;
\item notez présent ou absent, ainsi que la preuve exacte qui le montre ;
\item calculez vous-même le taux de couverture global ;
\item qualifiez ensuite la couverture zone par zone, pas seulement globalement ;
\item comparez seulement à ce stade avec le diagnostic JSON produit par le script.
\end{enumerate}
\end{frame}

\begin{frame}[fragile]{TP 3 -- Affichez d'abord les cinq attendus}
\small
\begin{verbatim}
head -n 6 data/samples/batch002_expected_sensors.csv
tail -n +2 data/samples/batch002_expected_sensors.csv | wc -l
\end{verbatim}
\begin{block}{Avant le calcul}
Qui autorise ce référentiel de cinq capteurs attendus ? Pour quel site, quelle mission et quelle date reste-t-il valide ?
\end{block}
\end{frame}

\begin{frame}[fragile]{TP 3 -- Contrôlez votre matrice avec le diagnostic}
\small
\begin{verbatim}
python3 -m json.tool data/processed/batch002_completeness.json
\end{verbatim}
Résultat attendu : 4 topics observés sur 5 attendus, \texttt{complete=False}, confiance faible, et le topic optronique explicitement signalé comme absent.
\begin{alertblock}{Formulation exacte à employer}
« 80 \% des topics attendus ont été observés », et jamais « 80 \% de la base est en sécurité ».
\end{alertblock}
\end{frame}

\begin{frame}{TP 3 -- Six questions pour qualifier le manque}
\small
\begin{enumerate}
\item Que signifie exactement le rapport 4/5 dans ce contexte ?
\item Pourquoi 80 \% ne mesure-t-il pas le risque réellement couvert par la base ?
\item Quelle zone présente une couverture strictement nulle ?
\item L'absence de cette zone prouve-t-elle une panne, ou pourrait-elle correspondre à une température normale ?
\item Quel test permettrait de départager un filtre erroné d'une réelle absence dans le broker ?
\item Quel contrôle renseignerait sur l'état physique réel de la zone, indépendamment des données ?
\end{enumerate}
\end{frame}

\begin{frame}{Le résultat numérique est 4/5, mais la décision dépend de la zone}
\centering\Huge 80\,\%\\[.5em]
\normalsize Quatre topics attendus sur cinq sont effectivement observés. L'abri optronique, pourtant classé en criticité haute par le référentiel, n'est couvert par aucun message.
\begin{alertblock}{Conséquence}
80 \% de couverture globale ne signifie en aucun cas « la base est couverte à 80 \% du risque réel » : la zone manquante est justement l'une des plus critiques.
\end{alertblock}
\end{frame}

\begin{frame}{Revue croisée -- attaquez le diagnostic, pas les personnes}
Le décideur critique vérifie systématiquement : le filtre utilisé, l'instant d'extraction, la branche interrogée, l'unicité des lignes, l'autorité du référentiel, la criticité de chaque zone et la formulation exacte de la conclusion.
\begin{block}{Question rouge}
Quelle hypothèse, non encore testée, pourrait transformer une absence apparemment réelle en une simple absence artificielle liée à notre propre extraction ?
\end{block}
\end{frame}

\section{Incident}
\begin{frame}{Incident injecté -- le capteur optronique attendu est absent}
\large Aucun topic \texttt{optronics-shelter-01/temperature} n'apparaît dans votre inventaire, alors qu'il figure bien dans le référentiel attendu.
\decisionquestion{Que savons-nous avec certitude, que supposons-nous seulement, et quelle vérification permettrait de départager ces hypothèses ?}
\end{frame}

\begin{frame}{Le silence admet plusieurs causes concurrentes}
\begin{enumerate}
\item le capteur lui-même, ou sa passerelle de transmission, peut être tombé en panne ;
\item aucune publication retained n'a peut-être jamais été émise pour ce topic ;
\item un message retained existait mais a pu être supprimé depuis ;
\item le filtre ou la branche interrogée peut être incorrecte de notre côté ;
\item l'extraction elle-même a pu être interrompue avant de recevoir ce message ;
\item le référentiel attendu peut être obsolète et mentionner un capteur qui n'existe plus.
\end{enumerate}
Une cause plausible parmi ces six n'est en aucun cas une cause démontrée : chacune reste à vérifier séparément.
\end{frame}

\begin{frame}{Une bonne vérification discrimine au moins deux hypothèses}
\small
\begin{tabular}{@{}p{.42\textwidth}p{.5\textwidth}@{}}
\toprule
Vérification & Hypothèses séparées \\
\midrule
répéter avec une branche plus large & filtre incorrect vs absence réelle dans le broker \\[2pt]
journaux du publisher ou de la passerelle & absence de publication vs problème de transport \\[2pt]
contrôle de la configuration retain & retained jamais activé vs retained supprimé \\[2pt]
contact du responsable du parc capteurs & référentiel valide vs référentiel obsolète \\[2pt]
mesure terrain indépendante & silence dans les données vs état physique réel \\
\bottomrule
\end{tabular}
\end{frame}

\begin{frame}{Incident -- Documentez chaque hypothèse avant de choisir}
\small
Pour au moins quatre des six hypothèses précédentes, renseignez :
\begin{enumerate}
\item une observation qui serait compatible avec cette hypothèse ;
\item la vérification concrète qui permettrait de la tester ;
\item le résultat qui la renforcerait si vous l'obteniez ;
\item le résultat qui l'affaiblirait au contraire ;
\item une action provisoire et réversible en attendant la vérification.
\end{enumerate}
\decisionquestion{Quelle vérification réduit l'incertitude le plus rapidement, sans jamais transformer ce silence en une mesure inventée ?}
\end{frame}

\begin{frame}{L'absence ne vaut jamais une température normale}
\begin{alertblock}{Surconclusion interdite}
L'absence totale de message ne fournit strictement aucune valeur thermique pour l'abri optronique -- ni basse, ni normale, ni élevée.
\end{alertblock}
La décision prudente réduit le périmètre de la conclusion à ce qui est réellement couvert, diminue explicitement la confiance globale, et priorise une vérification ciblée avant toute affirmation sur cette zone.
\end{frame}

\begin{frame}{Auto-vérification avant de continuer}
\small
\begin{block}{Le silence d'un topic prouve-t-il qu'un capteur est en panne ?}
Non : c'est une hypothèse compatible parmi six, pas une conclusion. Filtre incorrect, absence de publication, retained supprimé, extraction interrompue ou référentiel obsolète produiraient exactement le même silence observé.
\end{block}
\begin{block}{Que peut-on affirmer avec certitude à partir de ce seul silence ?}
Seulement que ce topic précis n'apparaît pas dans cette observation précise, avec ce filtre précis. Rien de plus sur sa cause, ni sur l'état physique réel de la zone.
\end{block}
\begin{block}{Comment choisir la meilleure vérification à mener en premier ?}
En comparant le délai nécessaire, la valeur d'information apportée et le risque opérationnel : une mesure terrain prioritaire si l'enjeu matériel est immédiat, des journaux techniques en parallèle si leur coût est faible.
\end{block}
\end{frame}

\section{TP 4 -- Décider}
\begin{frame}{TP 4 -- Le brief tient en cinq éléments vérifiables}
\begin{enumerate}
\item une action et un périmètre géographique exact et assumé ;
\item une confiance explicitement justifiée ;
\item deux preuves chiffrées et retrouvables dans les fichiers produits ;
\item deux incertitudes qui pourraient, si elles se confirmaient, changer la décision ;
\item une vérification à la fois prioritaire et réellement faisable avant 10 h.
\end{enumerate}
120 mots maximum, pas un de plus.
\end{frame}

\begin{frame}{TP 4 -- Procédez en trois passes, 15 minutes}
\begin{enumerate}
\item \textbf{5 min :} rédiger le brief avec action, périmètre, confiance, preuves, incertitudes et vérification ;
\item \textbf{4 min :} répondre sans détour aux cinq questions du décideur critique ;
\item \textbf{6 min :} réviser le texte, citer précisément les fichiers, voter et expliquer tout écart avec le vote initial.
\end{enumerate}
\begin{block}{Critère de fin}
Aucune phrase du brief ne généralise les quatre zones réellement observées à l'ensemble de la base.
\end{block}
\end{frame}

\begin{frame}{TP 4 -- Le décideur doit obtenir cinq réponses précises}
\small
\begin{enumerate}
\item Votre recommandation vaut-elle pour les quatre zones observées, ou pour les cinq zones de la base entière ?
\item Dans quel fichier exact retrouve-t-on chacune de vos deux preuves ?
\item Quelle hypothèse, si elle se vérifiait, pourrait renverser votre décision ?
\item Que proposez-vous concrètement pour l'abri optronique avant 10 h ?
\item Votre action reste-t-elle réversible si le diagnostic venait à changer ?
\end{enumerate}
\end{frame}

\begin{frame}{Une recommandation recevable sépare analyse partielle et décision globale}
\begin{block}{Structure possible}
Poursuivre l'analyse pour les quatre zones effectivement observées, sans généraliser cette analyse à toute la base. Confiance faible pour toute conclusion globale. Vérifier en priorité la chaîne optronique, ou effectuer une mesure terrain si le délai le permet.
\end{block}
À vous d'y insérer vos preuves et vos incertitudes exactes, chiffrées et retrouvables.
\end{frame}

\begin{frame}{Le décideur critique teste la portée de chaque phrase}
\begin{itemize}
\item « Pour quelles zones précises cette phrase est-elle vraie ? »
\item « Dans quel fichier puis-je retrouver cette preuve exacte ? »
\item « Quelle hypothèse renverserait entièrement l'action proposée ? »
\item « Quelle vérification reste possible avant l'échéance de 10 h ? »
\item « L'action proposée est-elle réversible si l'on se trompe ? »
\end{itemize}
\end{frame}

\begin{frame}{Vote final -- votre choix ou votre confiance ont-ils changé ?}
A couverture suffisante \quad B inspection ciblée \quad C suspendre la conclusion \quad D impossible sans inventaire
\begin{block}{À comparer explicitement au vote initial}
Le choix a-t-il changé ? La confiance a-t-elle changé ? Quelle preuve nouvelle, et quelle limite mieux comprise, expliquent cet écart ?
\end{block}
Une baisse de confiance par rapport au vote initial peut très bien signaler un meilleur diagnostic, pas un échec de raisonnement.
\end{frame}

\section{Synthèse}
\begin{frame}{Le réflexe à conserver tient en quatre verbes}
\centering\Large Observer $\rightarrow$ Inventorier $\rightarrow$ Comparer $\rightarrow$ Borner\\[1em]
\normalsize Observer avec un filtre explicite et documenté ; inventorier sans jamais perdre le contexte de réception ; comparer systématiquement à un attendu autorisé et daté ; borner enfin la décision au périmètre réellement démontré, ni plus, ni moins.
\end{frame}

\begin{frame}{Exit ticket -- trois phrases, sans écran}
\begin{enumerate}
\item Le broker permet d'affirmer que...
\item Il ne permet pas d'affirmer que...
\item Pour qualifier l'absence optronique, je vérifierais d'abord...
\end{enumerate}
\end{frame}

\begin{frame}{Réponse à la question directrice}
\decisionquestion{Avant tout nettoyage, peut-on savoir si le lot est suffisant pour soutenir une décision ?}
Oui, à condition de définir explicitement un périmètre et de comparer l'observé à un attendu tout aussi explicite. Ici, le lot n'est pas suffisant pour conclure sur l'ensemble de la base : 4 topics sur 5 attendus sont observés, et le topic manquant -- l'abri optronique -- est précisément l'un des plus critiques du référentiel.
\confidencelevel{faible pour toute décision couvrant la base entière ; moyenne, au mieux, pour décrire fidèlement l'instantané des quatre zones réellement observées.}
\end{frame}

\begin{frame}{Références}
\small
\begin{itemize}
\item OASIS Open, \emph{MQTT Version 5.0}, 2019. \url{https://docs.oasis-open.org/mqtt/mqtt/v5.0/os/mqtt-v5.0-os.html}
\item Eclipse Mosquitto, \emph{mosquitto\_sub manual}. \url{https://mosquitto.org/man/mosquitto_sub-1.html}
\item Eclipse Mosquitto, \emph{mosquitto\_pub manual}. \url{https://mosquitto.org/man/mosquitto_pub-1.html}
\end{itemize}
\end{frame}

\end{document}
